\section{Related Work}
\label{sec:related}

\subsection{Static Model Merging}
Model merging has emerged as a key strategy to combine specialized capabilities of multiple models fine-tuned from a shared progenitor. Pioneer works like Task Arithmetic \cite{Wortsman2022} demonstrated that adding task-specific fine-tuning weight vectors can produce multi-task models. However, simple linear blending often suffers from interference and parameter conflict. To alleviate this, TIES-Merging \cite{Yadav2023} resolves sign disagreements and prunes redundant parameters, while DARE \cite{Yu2023} employs random delta-weight dropping with rescaling to maintain representational stability. While highly effective, these approaches are static: they compute a single fixed set of interpolation coefficients during merging. Consequently, they cannot dynamically adjust expert blending based on the varying demands of individual incoming samples at runtime.

\subsection{Dynamic Model Merging and Routing}
To overcome the limitations of static merging, test-time dynamic model merging (or dynamic routing) was developed. Inspired by Mixture-of-Experts (MoE) architectures \cite{Shazeer2017, DudaHart2nd}, dynamic merging uses an input-dependent routing module to predict expert weights for each batch. This is highly practical when combined with parameter-efficient fine-tuning (PEFT) methods, where lightweight LoRA experts are dynamically interpolated. However, training standard parametric routers requires a calibration dataset representing all target tasks. When this calibration set is limited, standard linear routers suffer from transductive overfitting and representation collapse.

\subsection{Regularization and Parameter-Free Fallbacks}
Recent efforts have targeted the stability of dynamic routers. Task-Variance Regularization (VR-Router) penalizes the variance of routing coefficients to avoid vectorization collapse. Task-Space Anchor Regularization (TSAR) anchors routing weights directly to pre-calculated task-space centroids, restricting the router's drift under small-data training. On the other extreme, Parameter-Free Subspace Routing (PFSR) bypasses training altogether. PFSR utilizes a training-free projection technique where penultimate representations are projected onto pre-trained expert classification weight manifolds using cosine similarities. While PFSR is completely immune to overfitting, it lacks the adaptive, fine-grained calibration capability that a trained parametric router can acquire when sufficient data is present.

\subsection{The Empirical Gap}
Despite the proliferation of regularized and parameter-free routers, prior work lacks a systematic, rigorous, and stress-tested validation of these frameworks under realistic, non-idealized deployment scenarios. Specifically, two key questions remain unaddressed in existing literature:
\begin{enumerate}
    \item \emph{How do different routing mechanisms scale in terms of sample complexity, especially in the small-data regime ($N \le 32$)?}
    \item \emph{What happens to dynamic routing dynamics when deploying models on heterogeneous, mixed-task streams rather than clean, homogeneous batches?}
\end{enumerate}
In this work, we adopt a strictly empirical approach to fill these gaps. Through massive parallel sweeps over 5 independent seeds, we systematically map the transition between parametric and non-parametric regimes, and propose a complete, robust system architecture (CGHR + MBH) designed and verified specifically for deployment-level resilience.
